\documentclass[11pt,a4paper]{article}
\usepackage[utf8]{inputenc}
\usepackage[T2A]{fontenc}
\usepackage[english,russian]{babel}
\usepackage{amsmath,amssymb,amsthm}
\usepackage{algorithm}
\usepackage{algpseudocode}
\usepackage{geometry}
\usepackage{hyperref}
\usepackage{booktabs}
\usepackage{microtype}

\geometry{margin=1in}

\title{Иерархическая устойчивость GRA-обнулёнки: вариационная формулировка}
\author{bitsoev oleg}
\date{Май 2026}

\newtheorem{theorem}{Теорема}
\newtheorem{lemma}[theorem]{Лемма}
\newtheorem{axiom}{Аксиома}
\newtheorem{definition}{Определение}

\begin{document}

\selectlanguage{russian}

\maketitle

\begin{abstract}
Предлагается формализация механизма GRA-обнулёнки в виде вариационного иерархического функционала, определяющего динамику «пены» $\Phi$ в многоуровневой системе. Производная $d\Phi/dh$ по параметру уровня $h$ позволяет различать продуктивное схлопывание противоречий, стагнацию и деградацию. Сформулированы аксиомы, определения, леммы и теорема об устойчивости обнуления, а также дискретный алгоритм с псевдокодом для практической реализации.
\end{abstract}

\section{Введение}

Фреймворк GRA (Gradient Reduction of Argumentative Foam) предназначен для минимизации «пены» $\Phi$ --- меры противоречий, шума и избыточности в иерархических системах. Базовый механизм обнулёнки предполагает, что верхний уровень иерархии стремится к $\Phi \approx 0$, что интерпретируется как глобальное согласование. В данной работе обнулёнка рассматривается не только статически, но и динамически --- через производную $d\Phi/dh$. Это позволяет ввести критерий устойчивости обнуления и исследовать чувствительность системы к изменениям иерархии.

\section{Иерархическая модель и функционал пены}

Рассматривается система с уровнями $l = 0,\dots,L$. Состояние уровня --- $x_l \in \mathbb{R}^{n_l}$. Верхний уровень $L$ задаёт глобальное согласование. Локальная пена: $\Phi_l = \Phi(x_l) \geq 0$.

\textbf{Иерархический GRA-функционал:}
\begin{equation}
J[x_0,\dots,x_L] = \sum_{l=0}^L \alpha_l \Phi(x_l) + \sum_{l=0}^{L-1} \beta_l C_l(x_l, A_l x_{l+1}) + \gamma \Psi(x_L),
\label{eq:gra-functional}
\end{equation}
где:
\begin{itemize}
    \item $C_l$ --- штраф за несогласование между соседними уровнями,
    \item $A_l$ --- оператор подъёма/проекции,
    \item $\Psi(x_L)$ --- верхнеуровневое условие обнуления,
    \item $\alpha_l, \beta_l, \gamma > 0$ --- весовые коэффициенты.
\end{itemize}

При непрерывной параметризации уровня $h$, $\Phi(h)$ дифференцируема, и $d\Phi/dh$ --- градиент нулевости.

\section{Аксиомы GRA-обнулёнки}

\begin{axiom}[Иерархичность]
Система организована в уровни $l = 0,\dots,L$, каждый со своим состоянием.
\end{axiom}

\begin{axiom}[Пена]
Для каждого уровня определён функционал пены $\Phi(x_l) \geq 0$.
\end{axiom}

\begin{axiom}[Обнуление]
Существует операция $T_l$, уменьшающая пену: $\Phi(T_l(x_l)) \leq \Phi(x_l)$.
\end{axiom}

\begin{axiom}[Монотонность]
В идеальном режиме $\Phi_{l+1} \leq \Phi_l$.
\end{axiom}

\begin{axiom}[Чувствительность]
$\Phi(h)$ дифференцируема, знак $d\Phi/dh$ определяет режим эволюции.
\end{axiom}

\section{Определения}

\begin{definition}[GRA-функционал]
См. формулу~\ref{eq:gra-functional}.
\end{definition}

\begin{definition}[Устойчивое обнуление]
В точке $h^*$: $\Phi(h^*) = 0$, $\frac{d\Phi}{dh}(h^*) = 0$, $\frac{d^2\Phi}{dh^2}(h^*) > 0$.
\end{definition}

\begin{definition}[Дискретная производная пены]
$\Delta\Phi_l = \Phi_{l+1} - \Phi_l$, её дискретный аналог $d\Phi/dh$.
\end{definition}

\begin{definition}[Градиент нулевости]
$\nabla_h \Phi := \frac{d\Phi}{dh}$, характеризует направление и скорость изменения пены вдоль иерархии.
\end{definition}

\section{Леммы}

\begin{lemma}[Локальная стационарность]
В минимуме вариационные производные по каждому $x_l$ равны нулю.
\end{lemma}

\begin{lemma}[Нулевая пена не гарантирует устойчивости]
Необходимо также $\frac{d\Phi}{dh} = 0$.
\end{lemma}

\begin{lemma}[Положительная производная --- деградация]
Если $\frac{d\Phi}{dh} > 0$, пена растёт.
\end{lemma}

\begin{lemma}[Дискретный критерий устойчивости]
Для уровня $l^*$: $\Phi_{l^*} = 0$, $\Phi_{l^*+1} - \Phi_{l^*} = 0$, $\Phi_{l^*+1} - 2\Phi_{l^*} + \Phi_{l^*-1} > 0$.
\end{lemma}

\section{Теорема об устойчивости обнуления}

\begin{theorem}
Пусть $\Phi(h)$ дважды дифференцируема по $h$ и существует точка $h^*$ такая, что
\begin{equation}
\Phi(h^*) = 0, \quad \frac{d\Phi}{dh}(h^*) = 0, \quad \frac{d^2\Phi}{dh^2}(h^*) > 0.
\label{eq:stability-conditions}
\end{equation}
Тогда обнуление на уровне $h^*$ устойчиво: в некоторой окрестности возмущения не приводят к росту пены, и система возвращается к $\Phi \approx 0$.
\end{theorem}

\section{Набросок доказательства}

\begin{enumerate}
    \item Из $\frac{d\Phi}{dh}(h^*) = 0$ следует, что $h^*$ --- критическая точка.
    \item Условие $\frac{d^2\Phi}{dh^2}(h^*) > 0$ означает локальный минимум.
    \item Существует окрестность $U(h^*)$, где $\Phi(h) \geq \Phi(h^*) = 0$.
    \item Так как $\Phi \geq 0$, нулевая пена является локально минимальной.
    \item При малых возмущениях градиентный процесс возвращает систему к $h^*$, что и означает устойчивость. В дискретном случае рассуждение повторяется через конечные разности.
\end{enumerate}

\section{Дискретный алгоритм и псевдокод}

Алгоритм совмещает GRA-обнуление, выравнивание градиента $d\Phi/dh$ и контроль второй производной.

\begin{algorithm}[H]
\caption{HierarchicalStabilityOptimizer}
\begin{algorithmic}[1]
\State \textbf{class} HierarchicalStabilityOptimizer
\State \quad \textbf{def} \_\_init\_\_(self, hierarchy, lr=0.01):
\State \quad \quad self.h = hierarchy \Comment{список уровней с .state и .foam()}
\State \quad \quad self.lr = lr
\State \quad \textbf{def} step(self):
\State \quad \quad \# 1. GRA-обнуление
\State \quad \quad \textbf{for} level \textbf{in} self.h: level.state = gra\_nullify(level.state)
\State \quad \quad \# 2. Выравнивание dPhi/dh
\State \quad \quad \textbf{for} l \textbf{in} range(len(self.h)-1):
\State \quad \quad \quad dPhi = self.h[l+1].foam() - self.h[l].foam()
\State \quad \quad \quad \textbf{if} abs(dPhi) > 1e-6:
\State \quad \quad \quad \quad grad = grad\_connection(self.h[l], self.h[l+1], dPhi)
\State \quad \quad \quad \quad self.h[l+1].state -= self.lr * grad
\State \quad \quad \# 3. Обеспечение d\^2Phi/dh\^2 > 0
\State \quad \quad \textbf{for} l \textbf{in} range(1, len(self.h)-1):
\State \quad \quad \quad d2Phi = self.h[l+1].foam() - 2*self.h[l].foam() + self.h[l-1].foam()
\State \quad \quad \quad \textbf{if} d2Phi <= 0:
\State \quad \quad \quad \quad self.h[l].state += self.lr * 0.01 * (self.h[l+1].state + self.h[l-1].state - 2*self.h[l].state)
\State \quad \textbf{def} optimize(self, max\_iter=100):
\State \quad \quad \textbf{for} \_ \textbf{in} range(max\_iter): self.step()
\end{algorithmic}
\end{algorithm}

Полная реализация с абстрактными функциями находится в \texttt{src/optimizer.py}.

\section{Обсуждение и применения}

Интерпретация знака производной:
\begin{itemize}
    \item $\frac{d\Phi}{dh} < 0$ --- продуктивное схлопывание пены
    \item $\frac{d\Phi}{dh} \approx 0$ --- стагнация или плато
    \item $\frac{d\Phi}{dh} > 0$ --- деградация (галлюцинации, рост противоречий)
\end{itemize}

\textbf{Применения:}
\begin{itemize}
    \item Диагностика когнитивной устойчивости ИИ-архитектур (мониторинг «пены» в скрытых состояниях нейронных сетей)
    \item Многоуровневое управление в организациях (анализ согласованности решений на разных уровнях иерархии)
    \item Анализ эволюции научных теорий (отслеживание противоречий в системе знаний)
\end{itemize}

\section{Установка и использование}

\begin{verbatim}
git clone https://github.com/qqewq/GRA-Hierarchical-Stability.git
cd GRA-Hierarchical-Stability
pip install -r requirements.txt
python examples/demo.py
\end{verbatim}

Для интеграции с GRA-Multiverse-Final замените заглушки в \texttt{src/foam\_utils.py} на реальные функции из вашего проекта.

\section{Ссылки}

\begin{itemize}
    \item GRA-Multiverse-Final
    \item GRA-Swarm
    \item NeurIPS 2022 hierarchical optimization materials
    \item HAL preprint hal-04088837
    \item SPP1962 preprint 036, WIAS Berlin
\end{itemize}

\end{document}